\section{Theory}

This section is currently being developed alongside the implementation. The immediate objective is to establish a verified three-dimensional finite element benchmark for compression of a single elastic cylinder between frictionless rigid platens. The benchmark will provide the first working mechanics problem for the project and establish the computational patterns that will later be extended to granular assemblies.

The implementation roadmap is listed below. Each item identifies the principal source files expected to implement that capability. File names may change as the implementation develops; the list is intended to make the dependency structure of the first simulation explicit.

\subsection{Implementation Roadmap}

\subsection{Initial Benchmark: 3D Linear Elastic Cylinder}

The initial benchmark is a three-dimensional, displacement-controlled
compression test of a linear-elastic cylinder. The benchmark is developed
incrementally so that the physical model, numerical formulation, solution
procedure, and validation remain separate and independently testable.

The current implementation follows the modular architecture

$$
\texttt{problem}
\;\rightarrow\;
\texttt{mechanics}
\;\rightarrow\;
\texttt{numerical}
\;\rightarrow\;
\texttt{io}.
$$

Within the numerical layer, the current steady-state solution path is

$$
\texttt{formulation}
\;\rightarrow\;
\texttt{boundary}
\;\rightarrow\;
\texttt{solver}.
$$

The current development status is:

\begin{longtable}{p{0.06\textwidth}p{0.36\textwidth}p{0.32\textwidth}p{0.16\textwidth}}
\toprule
\textbf{ID} & \textbf{Task} & \textbf{Relevant files} & \textbf{Status} \\
\midrule
\endhead

1 &
Define the 3D cylinder geometry &
\texttt{problem/geometry/geometry\_types/cylinder.py} &
Implemented \\

2 &
Generate the tetrahedral finite-element mesh and expose named physical regions &
\texttt{problem/geometry/meshing.py},
\texttt{problem/geometry/geometry\_types/cylinder.py} &
Implemented \\

3 &
Define small-strain kinematics &
\texttt{mechanics/kinematics/small\_strain.py} &
Implemented \\

4 &
Implement isotropic linear elasticity &
\texttt{mechanics/constitutive/linear\_elastic.py} &
Implemented \\

5 &
Define the compression boundary constraints and rigid-body constraints &
\texttt{problem/definition.py},
\texttt{problem/loader.py},
\texttt{mechanics/rigid\_body.py} &
Implemented \\

6 &
Translate problem-level Dirichlet constraints into DOLFINx boundary conditions &
\texttt{numerical/boundary/dirichlet.py} &
Implemented \\

7 &
Construct the finite-element elasticity formulation from the mechanics models &
\texttt{numerical/formulation/elasticity.py} &
Implemented \\

8 &
Construct the reference measure and rigid-body constraint functionals &
\texttt{numerical/boundary/reference.py},
\texttt{numerical/formulation/rigid\_body.py} &
Implemented \\

9 &
Solve the steady linear-elasticity system with DOLFINx/PETSc &
\texttt{numerical/solver/steady.py} &
Implemented \\

10 &
Run the complete small-strain cylinder compression problem end-to-end &
\texttt{numerical/},
\texttt{\_\_main\_\_.py},
\texttt{tests/integration/} &
Implemented \\

11 &
Provide command-line diagnostic reporting for the application execution path &
\texttt{\_\_main\_\_.py} &
Implemented \\

12 &
Extract reaction forces and macroscopic compression response &
\texttt{numerical/},
\texttt{io/} &
TODO \\

13 &
Write VTK-compatible simulation output &
\texttt{io/} &
TODO \\

14 &
Implement an independent analytical or semi-analytical reference solution &
\texttt{benchmarks/01\_cylinder\_compression/reference.py} &
TODO \\

15 &
Compare the finite-element solution with the reference solution &
\texttt{benchmarks/01\_cylinder\_compression/},
\texttt{tests/} &
TODO \\

16 &
Perform a mesh-convergence study &
\texttt{benchmarks/01\_cylinder\_compression/experiment.yaml},
\texttt{tests/} &
TODO \\

17 &
Implement frictionless platen/contact treatment &
\texttt{mechanics/contact/},
\texttt{problem/},
future numerical contact layer &
TODO \\

18 &
Add finite-strain kinematics &
\texttt{mechanics/kinematics/finite\_strain.py} &
TODO \\

19 &
Implement a compressible Neo-Hookean constitutive model &
\texttt{mechanics/constitutive/neo\_hookean.py} &
TODO \\

20 &
Add the nonlinear finite-strain solver &
\texttt{numerical/} &
TODO \\

21 &
Compare small-strain and finite-strain solutions &
\texttt{benchmarks/01\_cylinder\_compression/} &
TODO \\

22 &
Compute the relative $L^2$ strain-field error &
\texttt{benchmarks/01\_cylinder\_compression/},
\texttt{io/} &
TODO \\

\bottomrule
\end{longtable}

The immediate milestone is now to validate the completed small-strain
linear-elastic solution path against independent physical and numerical
checks. The implemented computational path is

$$
\text{ProblemDefinition}
\rightarrow
\text{GeometryDefinition}
\rightarrow
\text{MeshData}
\rightarrow
\text{ElasticityFormulation}
\rightarrow
\text{Dirichlet BCs}
\rightarrow
\text{Rigid-body constraints}
\rightarrow
\text{SteadySolver}.
$$

The current implementation has therefore completed the basic
end-to-end solution stage. The next development stages are to expose
the computed response, write simulation output, and establish independent
analytical and mesh-convergence validation before extending the model to
contact and finite-strain mechanics.

The purpose of this problem is not to represent granular mechanics yet. It is a controlled computational mechanics problem that allows the basic FEniCSx implementation, boundary treatment, solver configuration, output, and verification workflow to be established before introducing multiple grains, evolving contact networks, periodicity, and realistic grain geometries.

\subsection{Reference Solution and Verification}

The reference solution will be maintained independently of the core finite element implementation. In particular, the core package will not contain axisymmetric-specific logic merely to support this benchmark.

The reference problem will exploit the axisymmetric structure of the cylindrical compression problem. The resulting displacement and stress fields are generally functions of radial and axial position and are not assumed to be spatially uniform near the platen interfaces. The reference calculation will therefore provide a physically meaningful comparison for the three-dimensional finite element solution rather than reducing the verification problem to a homogeneous uniaxial-stress calculation.

Verification quantities will include the macroscopic reaction force and, where appropriate, displacement and strain-field comparisons. Mesh refinement will be used to determine whether the three-dimensional finite element solution approaches the reference solution systematically.

\subsection{Reaction-force Definition and Evaluation}

The macroscopic reaction force is defined on the actual boundary of the finite-element domain on which the displacement constraint is imposed. No extrapolation of the stress field beyond the finite-element mesh is used to define the reaction.

For a loaded boundary $\Gamma_L$ with outward unit normal $\mathbf{n}$, the corresponding continuum resultant is

$$
\mathbf{F}_R
=
\int_{\Gamma_L}
\boldsymbol{\sigma}\mathbf{n}\,d\Gamma.
$$

In the present displacement-controlled formulation, the primary numerical reaction will instead be evaluated from the discrete finite-element equilibrium system. For the linear system

$$
\mathbf{K}\mathbf{u}=\mathbf{f},
$$

the residual associated with the prescribed displacement degrees of freedom represents the reaction required to enforce the Dirichlet constraints. The resultant reaction is therefore obtained from the constrained degrees of freedom of the discrete problem rather than from extrapolation of interior quantities.

The boundary traction resultant provides an independent verification quantity. Where the finite-element stress field is sufficiently regular on the boundary, the traction-based evaluation

$$
\mathbf{F}_{\Gamma}
=
\int_{\Gamma_L}
\boldsymbol{\sigma}\mathbf{n}\,d\Gamma
$$

will be compared with the reaction obtained from the discrete constrained system. Agreement between these two evaluations provides an additional equilibrium check on the implementation and post-processing.

The reaction will be retained as a vector quantity,

$$
\mathbf{F}_R =
\begin{bmatrix}
F_x & F_y & F_z
\end{bmatrix}^{T},
$$

rather than being reduced immediately to a scalar magnitude. A scalar compressive force will subsequently be obtained by projection onto the prescribed loading direction, with the sign convention stated explicitly. This representation avoids introducing a hidden assumption about the loading direction and permits the same response definition to be reused for other loading orientations and future contact formulations.

This distinction between the physical boundary resultant and its discrete numerical evaluation will be maintained throughout the benchmark. The same reaction-force definition will subsequently be used when constructing the macroscopic force--displacement and stress--strain response.


\subsection{Finite-Strain Extension}

After the small-strain benchmark has been verified, the same geometric and loading configuration will be extended to finite-strain elasticity. The finite-strain formulation will introduce the deformation gradient and a nonlinear constitutive response, initially using a compressible Neo-Hookean model whose small-strain limit is consistent with the selected linear-elastic material parameters.

The finite-strain calculation will use the nonlinear solution infrastructure provided by DOLFINx and PETSc. The resulting solution will be compared with the small-strain calculation at increasing levels of compression.

The comparison will include:

\begin{itemize}
\item reaction force versus applied engineering strain;
\item selected stress and strain fields;
\item displacement fields; and
\item the relative $L^2$ error between the finite-strain and small-strain strain fields.
\end{itemize}

The strain-field comparison will be used to identify the deformation range over which the small-strain approximation remains adequate and the range in which finite-strain effects become significant. The interpretation of this comparison will distinguish kinematic effects from constitutive-model differences.

\subsection{Planned Benchmark Structure}

The first benchmark is expected to be organized independently from the core mechanics package:

\begin{verbatim}
benchmarks/
`-- 01_cylinder_compression/     |-- README.md     |-- experiment.yaml
    `-- reference.py
\end{verbatim}

The benchmark configuration will describe the physical experiment and numerical parameters, while the reusable implementation will remain in the \texttt{granular\_rves} package. This separation is intended to keep benchmark-specific verification logic out of the general mechanics infrastructure.

As the implementation matures, this section will be replaced progressively by the mathematical formulation, discretization details, solver configuration, verification methodology, and results associated with each completed milestone.

